\documentclass[10pt,a4paper]{amsart}
\usepackage{amsmath,amssymb,amsthm,graphicx}
\usepackage{verbatim}
\newenvironment{reference}{\comment}{\endcomment}
%\newenvironment{reference}{}{}
\newenvironment{slogan}{\comment}{\endcomment}
\newenvironment{history}{\comment}{\endcomment}

% For commutative diagrams we use Xy-pic
\usepackage[all]{xy}

% We use 2cell for 2-commutative diagrams.
\xyoption{2cell}
% Theorem environments.
%
\theoremstyle{plain}
\newtheorem{theorem}[subsection]{Theorem}
\newtheorem{proposition}[subsection]{Proposition}
\newtheorem{lemma}[subsection]{Lemma}

\theoremstyle{definition}
\newtheorem{definition}[subsection]{Definition}
\newtheorem{example}[subsection]{Example}
\newtheorem{exercise}[subsection]{Exercise}
\newtheorem{situation}[subsection]{Situation}

\theoremstyle{remark}
\newtheorem{remark}[subsection]{Remark}
\newtheorem{remarks}[subsection]{Remarks}

\numberwithin{equation}{subsection}

% Macros
%
\def\lim{\mathop{\mathrm{lim}}\nolimits}
\def\colim{\mathop{\mathrm{colim}}\nolimits}
\def\Spec{\mathop{\mathrm{Spec}}}
\def\Hom{\mathop{\mathrm{Hom}}\nolimits}
\def\Ext{\mathop{\mathrm{Ext}}\nolimits}
\def\SheafHom{\mathop{\mathcal{H}\!\mathit{om}}\nolimits}
\def\SheafExt{\mathop{\mathcal{E}\!\mathit{xt}}\nolimits}
\def\Sch{\mathit{Sch}}
\def\Mor{\mathop{\mathrm{Mor}}\nolimits}
\def\Ob{\mathop{\mathrm{Ob}}\nolimits}
\def\Sh{\mathop{\mathit{Sh}}\nolimits}
\def\NL{\mathop{N\!L}\nolimits}
\def\CH{\mathop{\mathrm{CH}}\nolimits}
\def\proetale{{pro\text{-}\acute{e}tale}}
\def\etale{{\acute{e}tale}}
\def\QCoh{\mathit{QCoh}}
\def\Ker{\mathop{\mathrm{Ker}}}
\def\Im{\mathop{\mathrm{Im}}}
\def\Coker{\mathop{\mathrm{Coker}}}
\def\Coim{\mathop{\mathrm{Coim}}}

% Boxtimes
%
\DeclareMathSymbol{\boxtimes}{\mathbin}{AMSa}{"02}

%
% Macros for moduli stacks/spaces
%
\def\QCohstack{\mathcal{QC}\!\mathit{oh}}
\def\Cohstack{\mathcal{C}\!\mathit{oh}}
\def\Spacesstack{\mathcal{S}\!\mathit{paces}}
\def\Quotfunctor{\mathrm{Quot}}
\def\Hilbfunctor{\mathrm{Hilb}}
\def\Curvesstack{\mathcal{C}\!\mathit{urves}}
\def\Polarizedstack{\mathcal{P}\!\mathit{olarized}}
\def\Complexesstack{\mathcal{C}\!\mathit{omplexes}}
% \Pic is the operator that assigns to X its picard group, usage \Pic(X)
% \Picardstack_{X/B} denotes the Picard stack of X over B
% \Picardfunctor_{X/B} denotes the Picard functor of X over B
\def\Pic{\mathop{\mathrm{Pic}}\nolimits}
\def\Picardstack{\mathcal{P}\!\mathit{ic}}
\def\Picardfunctor{\mathrm{Pic}}
\def\Deformationcategory{\mathcal{D}\!\mathit{ef}}
\usepackage{cleveref}
\crefname{originallemma}{original lemma}{original lemmas}
\Crefname{originallemma}{Original lemma}{Original lemmas}
\makeatletter
\newlabel{lemma-gysin-at-infty}{{0F9J}{source}{}{sourcepointer-1}{}}
\newlabel{lemma-gysin-at-infty@cref}{{[originallemma][0][]0F9J}{[1][0][]source}{}{}{}}
\makeatother
\crefname{originallemma}{original lemma}{original lemmas}
\Crefname{originallemma}{Original lemma}{Original lemmas}
\makeatletter
\newlabel{lemma-support-cap-effective-Cartier}{{02T9}{source}{}{sourcepointer-2}{}}
\newlabel{lemma-support-cap-effective-Cartier@cref}{{[originallemma][0][]02T9}{[1][0][]source}{}{}{}}
\makeatother
\crefname{originallemma}{original lemma}{original lemmas}
\Crefname{originallemma}{Original lemma}{Original lemmas}
\makeatletter
\newlabel{lemma-linebundle}{{0B74}{source}{}{sourcepointer-3}{}}
\newlabel{lemma-linebundle@cref}{{[originallemma][0][]0B74}{[1][0][]source}{}{}{}}
\makeatother
\crefname{originallemma}{original lemma}{original lemmas}
\Crefname{originallemma}{Original lemma}{Original lemmas}
\makeatletter
\newlabel{lemma-linebundle-formulae}{{0F96}{source}{}{sourcepointer-4}{}}
\newlabel{lemma-linebundle-formulae@cref}{{[originallemma][0][]0F96}{[1][0][]source}{}{}{}}
\makeatother
\title{The composite of the Gysin-at-infinity class with the pushforward from infinity equals the Gysin pullback at infinity}
\author{The Stacks Project Authors}
\date{stacks-project:chow.tex:local-snapshot-20261012; GNU FDL 1.2 or later}
\begin{document}
\maketitle
\noindent\textit{Editorial note.} All delivered mathematical text is the author's own, copied from the pinned Stacks Project chapter file chow.tex; no statement or proof step was rewritten, reordered or removed.\par\medskip
\noindent\textit{Editorial note.} Preamble provenance: the pinned chapter file opens with an input of the shared preamble and carries no preamble text of its own, so the project's single Stacks-wide preamble staged at sources/stacks-chow/preamble.tex is excerpted verbatim as two spans, lines 14-18 (Xy-pic setup) and lines 154-217 (the native theorem declarations and mathematical macros). Nothing was added, removed or redefined; the wrapper class and the amsmath/amssymb/amsthm/graphicx base packages are supplied by the assembler.\par\medskip
\noindent\textit{Editorial note.} The counted statement and proof contain no cross reference, no citation and no included file, so the delivered body is self contained.\par\medskip
\noindent\textit{Editorial note.} Canonical location of the counted proposition: Stacks tag 0FAR (sources/stacks-chow/pinned-original-source.tex line 10132).\par\medskip
\noindent\textit{Allowed standard prerequisites (editorial).}\begin{itemize}
\item The Gysin-at-infinity class and its defining property. Same chapter, Lemma labelled lemma-gysin-at-infty (source line 10026), invoked in the statement. Supplies the class C and the characterisation C cap alpha = i\_infinity\textasciicircum{}* beta which is what the proof must match.
\item The two rationally equivalent effective Cartier divisors of the projective line give the same cycle class. Same chapter, Lemma labelled lemma-support-cap-effective-Cartier (source line 5138), invoked at proof lines 10142-10143. Supplies that i\_\{infinity,*\}i\_infinity\textasciicircum{}* beta and i\_\{0,*\}i\_0\textasciicircum{}* beta represent the same class on W.
\item The line bundle lemmas identifying fibre restrictions after pullback along the zero section. Same chapter, Lemmas labelled lemma-linebundle (source line 5737) and lemma-linebundle-formulae (source line 5775), invoked at proof lines 10148-10149. Supply that the restriction of beta to b\textasciicircum{}\{-1\}(A\textasciicircum{}1\_X) restricts to the flat pullback of i\_0\textasciicircum{}* beta, which may be verified after pullback along i\_0.
\end{itemize}\medskip
\ifdefined\hypertarget\hypertarget{sourcepointer-1}{}\fi
\noindent\textit{Source reference.} Original lemma 0F9J: sources/stacks-chow/pinned-original-source.tex, source lines 10026--10026. Same-chapter pointer into the author environment lemma at source line 10026; invoked at proof line None of the counted proof. It is retained only as a navigation pointer; every load-bearing step is inside the delivered span.\par\medskip
\ifdefined\hypertarget\hypertarget{sourcepointer-2}{}\fi
\noindent\textit{Source reference.} Original lemma 02T9: sources/stacks-chow/pinned-original-source.tex, source lines 5138--5138. Same-chapter pointer into the author environment lemma at source line 5138; invoked at proof line 10143 of the counted proof. It is retained only as a navigation pointer; every load-bearing step is inside the delivered span.\par\medskip
\ifdefined\hypertarget\hypertarget{sourcepointer-3}{}\fi
\noindent\textit{Source reference.} Original lemma 0B74: sources/stacks-chow/pinned-original-source.tex, source lines 5737--5737. Same-chapter pointer into the author environment lemma at source line 5737; invoked at proof line 10149 of the counted proof. It is retained only as a navigation pointer; every load-bearing step is inside the delivered span.\par\medskip
\ifdefined\hypertarget\hypertarget{sourcepointer-4}{}\fi
\noindent\textit{Source reference.} Original lemma 0F96: sources/stacks-chow/pinned-original-source.tex, source lines 5775--5775. Same-chapter pointer into the author environment lemma at source line 5775; invoked at proof line 10149 of the counted proof. It is retained only as a navigation pointer; every load-bearing step is inside the delivered span.\par\medskip
\begin{lemma}
\label{lemma-homomorphism-pre}
In Lemma \ref{lemma-gysin-at-infty} we have
$C \circ (W_\infty \to X)_* \circ i_\infty^* = i_\infty^*$.
\end{lemma}

\begin{proof}
Let $\beta \in \CH_{k + 1}(W)$. Denote $i_0 : X = X \times \{0\} \to W$
the closed immersion of the fibre over $0$ in $\mathbf{P}^1$. Then
$(W_\infty \to X)_* i_\infty^* \beta = i_0^*\beta$ in $\CH_k(X)$ because
$i_{\infty, *}i_\infty^*\beta$ and $i_{0, *}i_0^*\beta$
represent the same class on $W$ (for example by
Lemma \ref{lemma-support-cap-effective-Cartier})
and hence pushforward to the same class on $X$.
The restriction of $\beta$ to $b^{-1}(\mathbf{A}^1_X)$
restricts to the flat pullback of
$i_0^*\beta = (W_\infty \to X)_* i_\infty^* \beta$ because we can check
this after pullback by $i_0$, see
Lemmas \ref{lemma-linebundle} and \ref{lemma-linebundle-formulae}.
Hence we may use $\beta$ when computing the image of
$(W_\infty \to X)_*i_\infty^*\beta$ under $C$
and we get the desired result.
\end{proof}

\end{document}
